%!TEX root = ../main.tex

\chapter{Cheat sheet}
\label{chp:cheat}


% -----------------------------


Random citation \cite{alecci2021development}.


% -----------------------------


Random footnote.\footnote{\url{https://lucamartinelli.eu.org}}


% -----------------------------


\section{A section}


% -----------------------------


\subsubsection{Example of list}
\begin{itemize}
  \item Item 1
  \item Item 2
\end{itemize}


% -----------------------------


\subsection{A subsection}


% -----------------------------


\subsubsection{Example of acronym}
\ac{NLP}


% -----------------------------


\subsubsection{Example of enumeration}
\begin{enumerate}
  \item Item 1
  \item Item 2
\end{enumerate}


% -----------------------------


\subsubsection{Example of quote}
\begin{displayquote}
Lorem ipsum dolor sit amet, consectetur adipiscing elit, sed do eiusmod tempor incididunt ut labore et dolore magna aliqua.
\end{displayquote}


% -----------------------------


\begin{figure}[ht]
    \centering
    \includegraphics[width=0.5\textwidth]{res/ltunipd}
    \caption{Example of image}
\end{figure}


% -----------------------------


\begin{algorithm}[ht]
    \caption{An algorithm with caption}\label{alg:two}
    \begin{algorithmic}
        \REQUIRE $n \geq 0$
        \ENSURE $y = x^n$
        
        \STATE $y \gets 1$
        \STATE $X \gets x$
        \STATE $N \gets n$
        
        \WHILE{$N \neq 0$}
            \IF{$N$ is even}
              \STATE $X \gets X \times X$
              \STATE $N \gets \frac{N}{2} $  \COMMENT{This is a comment}
            \ELSIF{$N$ is odd}
              \STATE $y \gets y \times X$
              \STATE $N \gets N - 1$
            \ENDIF
        \ENDWHILE
    \end{algorithmic}
\end{algorithm}


% -----------------------------


\begin{equation}
e^{j\pi} + 1 = 0
\end{equation}


% -----------------------------


\tikzset{vertex style/.style={
    draw=#1,
    thick,
    fill=#1!70,
    text=white,
    ellipse,
    minimum width=2cm,
    minimum height=0.75cm,
    font=\small,
    outer sep=3pt,
  },
  text style/.style={
    sloped,
    text=black,
    font=\footnotesize,
    above
  }
}


% -----------------------------


\begin{figure}[ht]
    \centering
    \begin{tikzpicture}[node distance=2.75cm,>={Stealth[]}]
        \node[vertex style=cyan] (Rk) {Righteous Kill};
        \node[vertex style=orange, above of=Rk,xshift=2em] (BD) {Bryan Dennehy}
        edge [<-,cyan!60!blue] node[text style,above]{starring} (Rk);
    \end{tikzpicture}
    
    \caption{Image created with TikZ} \label{fig:T1}
\end{figure}


% -----------------------------


\paragraph{}
Lorem ipsum dolor sit amet, consectetur adipiscing elit, sed do eiusmod tempor incididunt ut labore et dolore magna aliqua. Ut enim ad minim veniam, quis nostrud exercitation ullamco laboris nisi ut aliquip ex ea commodo consequat. Duis aute irure dolor in reprehenderit in voluptate velit esse cillum dolore eu fugiat nulla pariatur. Excepteur sint occaecat cupidatat non proident, sunt in culpa qui officia deserunt mollit anim id est laborum.


% -----------------------------


\begin{lstlisting}[language=Python, caption=Code snippet example]
import numpy as np
    
def incmatrix(genl1,genl2):
    m = len(genl1)
    n = len(genl2)
    M = None #to become the incidence matrix
    VT = np.zeros((n*m,1), int)  #dummy variable

    test = "String"
    
    #compute the bitwise xor matrix
    M1 = bitxormatrix(genl1)
    M2 = np.triu(bitxormatrix(genl2),1) 

    for i in range(m-1):
        for j in range(i+1, m):
            [r,c] = np.where(M2 == M1[i,j])
            for k in range(len(r)):
                VT[(i)*n + r[k]] = 1;
                VT[(i)*n + c[k]] = 1;
                VT[(j)*n + r[k]] = 1;
                VT[(j)*n + c[k]] = 1;
                
                if M is None:
                    M = np.copy(VT)
                else:
                    M = np.concatenate((M, VT), 1)
                
                VT = np.zeros((n*m,1), int)
    
    return M
\end{lstlisting}


% -----------------------------


\begin{table}[ht]
    \centering
    \begin{tabular}{|l|l|}
    \hline
    \textbf{A} & \textbf{B} \\ \hline
    C          & D          \\ \hline
    E          & F          \\ \hline
    G          & H          \\ \hline
    \end{tabular}
    \caption{Table example \label{tab:table-name}}
\end{table}